O algoritmo Bubble Sort apresenta comportamento semelhante nos diferentes
tipos de entrada, pois realiza comparações entre elementos adjacentes
durante sucessivas passagens pelo vetor. Sua complexidade de tempo é
analisada considerando o melhor caso, o caso médio e o pior caso.

\subsubsection{Melhor caso -- \texorpdfstring{$\Theta(n^2)$}{Theta(n²)}}

O melhor caso ocorre quando o vetor já está ordenado de forma crescente.
Mesmo nessa situação, a implementação utilizada continua realizando todas
as comparações previstas pelos laços do algoritmo, embora nenhuma troca
seja necessária.

Considerando um vetor com $n$ elementos, na primeira passagem são
realizadas $n-1$ comparações. Na segunda passagem, são realizadas $n-2$
comparações, e assim sucessivamente até a última passagem, que realiza
apenas uma comparação.

Assim, o número total de comparações pode ser representado pela soma:

\begin{equation}
T(n) = (n-1) + (n-2) + \cdots + 2 + 1
\end{equation}

Essa soma corresponde à soma dos primeiros $n-1$ números naturais:

\begin{equation}
T(n) = \frac{n(n-1)}{2}
\end{equation}

Desenvolvendo a expressão:

\begin{equation}
T(n) = \frac{n^2-n}{2}
\end{equation}

O termo de maior crescimento é $n^2$. Portanto:

\begin{equation}
T(n) = \Theta(n^2)
\end{equation}

Assim, o melhor caso do Bubble Sort apresenta complexidade quadrática.

\subsubsection{Pior caso -- \texorpdfstring{$\Theta(n^2)$}{Theta(n2)}}

O pior caso ocorre quando o vetor está ordenado de forma decrescente e
deseja-se ordená-lo de forma crescente. Nesse cenário, além de realizar
todas as comparações previstas pelos laços, o algoritmo também precisa
realizar uma grande quantidade de trocas entre os elementos.

A quantidade de comparações realizadas permanece igual à do melhor caso,
sendo representada por:

\begin{equation}
T(n) = (n-1) + (n-2) + \cdots + 2 + 1
\end{equation}

Aplicando a fórmula da soma:

\begin{equation}
T(n) = \frac{n(n-1)}{2}
\end{equation}

\begin{equation}
T(n) = \frac{n^2-n}{2}
\end{equation}

Além das comparações, são realizadas trocas dos elementos que estão fora
de ordem. No pior caso, a quantidade de trocas também cresce
quadraticamente.

Como o termo dominante é $n^2$, temos:

\begin{equation}
T(n) = \Theta(n^2)
\end{equation}

Portanto, o pior caso do Bubble Sort apresenta complexidade quadrática.

\subsubsection{Caso médio -- \texorpdfstring{$\Theta(n^2)$}{Theta(n²)}}

O caso médio ocorre quando os elementos estão distribuídos de maneira
aleatória. Nessa situação, o algoritmo continua realizando todas as
passagens e comparações previstas pelos seus laços.

A quantidade de comparações é dada por:

\begin{equation}
T(n) = (n-1) + (n-2) + \cdots + 2 + 1
\end{equation}

Aplicando a fórmula da soma:

\begin{equation}
T(n) = \frac{n(n-1)}{2}
\end{equation}

\begin{equation}
T(n) = \frac{n^2-n}{2}
\end{equation}

Embora a quantidade de trocas varie de acordo com a organização inicial
dos elementos, o número de comparações continua sendo proporcional a
$n^2$.

Como o termo dominante é $n^2$, a complexidade assintótica é:

\begin{equation}
T(n) = \Theta(n^2)
\end{equation}

Portanto, no caso médio, o Bubble Sort também apresenta complexidade
quadrática.

\subsubsection{Resumo das complexidades}

\begin{table}[H]
    \centering
    \caption{Complexidades do algoritmo Bubble Sort}
    \label{tab:complexidade-bubble-analise}
    \begin{tabular}{lc}
        \hline
        \textbf{Caso} & \textbf{Complexidade de tempo} \\
        \hline
        Melhor caso & $\Theta(n)$ \\
        Caso médio & $\Theta(n^2)$ \\
        Pior caso & $\Theta(n^2)$ \\
        \hline
    \end{tabular}
\end{table}

Conclui-se que o Bubble Sort apresenta complexidade quadrática nos três
casos analisados. Mesmo quando o vetor já está ordenado, a implementação
utilizada continua realizando as comparações das passagens previstas.
Por esse motivo, o algoritmo pode apresentar desempenho limitado para
grandes conjuntos de dados.
